\section*{3. 自编码器与自监督学习（6 分）}

本题使用 CIFAR-10 数据集，比较自编码器（Auto-Encoder）和两种自监督学习的表示学习效果。
三种方法的编码器架构、参数量和输出表示维度尽量保持一致。
在 CIFAR-10 的训练集上分别训练 3 个模型，模型训练阶段不得使用图像的类别标签。

\begin{enumerate}
    \item \textbf{自编码器：}
    实现包含编码器和解码器的自编码器，将重构误差作为训练损失。
    报告测试集上的重构误差，并展示测试集中原图和重构图的示例。
    （1 分）

    \item \textbf{旋转角度预测：}
    将训练集中的每张图像分别随机旋转 $0^\circ, 90^\circ, 180^\circ, 270^\circ$，
    实现编码器和旋转角度分类器，训练模型预测旋转角度。
    将测试集中的每张图像分别随机旋转 $0^\circ, 90^\circ, 180^\circ, 270^\circ$，
    报告测试集上旋转预测准确率。
    （1 分）

    \item \textbf{掩码自编码器（MAE）：}
    将训练集中的图像划分为块（patch），每个块的分辨率为 $4 \times 4$。
    随机遮蔽 $75\%$ 比例的块，每张图片每个 epoch 中的随机化方式可以不同，
    将被遮蔽块上的重构 MSE 作为训练损失。
    对测试集中的每张图片随机遮蔽 $75\%$ 比例的块，产生一张测试样本，
    报告测试样本集上被遮蔽块上的重构 MSE，并展示原图、遮蔽图和重构图的示例。
    注意，MAE 编码器的实现需要支持变长的输入，
    即既支持遮蔽 $75\%$ 比例块的输入，也支持不做遮蔽的原始所有块的输入。
    （1 分）

    \item 分别使用以上自编码器、旋转角度预测和 MAE 预训练得到的编码器，
    在 CIFAR-10 测试集上评估分类效果：
    冻结编码器权重，在训练集上，以编码器的输出作为线性分类器的输入，
    用交叉熵损失训练分类器。
    比较三种表示学习方法在测试集上的分类准确率，
    用 t-SNE 降维方法可视化测试集上学到的特征，
    并分析三种表示学习方法学到的表示有何不同，
    哪种方法对分类任务更有帮助。
    （3 分）
\end{enumerate}